\chapter{Litigation inference games and issue-specific burdens}
\label{sec:p18-lis-burdens}

A party can be required to establish one issue while needing only to
raise doubt on another. The Litigation Inference System, LIS, separates
the identity of the party from its dialectical role at a particular
argument and subargument.
Prakken and Sartor's 2009 account \cite{PrakkenSartor2009Burdens}
explains the burden distinctions and their changing evidential setting.
Its precise changing-role game is referred to Prakken's 2001
presentation \cite{Prakken2001Defeasibility}.
The name LIS and its relation to the earlier inference system IS
are made explicit in the 2006 paper
\cite{PrakkenSartor2006Presumptions}.

This chapter gives the dialectical construction with a fully specified
argument interface, history-dependent burden input, move rules,
winning conditions, and proofs. It distinguishes the fixed-role
grounded game from the changing-role litigation game.
The burden domain is explicitly issue-relative; the source's
closure condition and its obstruction under unrestricted classical
closure are characterized below. This additional domain condition is part of the corrected
construction, not an assertion that the unrepaired condition works
for unrestricted classical formulas.

LIS's argument game evaluates a supplied information state. It is not
by itself a complete legal procedure for admission, pleading,
concession, withdrawal, or judicial allocation of burdens.
The 2009 paper expressly abstracts from those procedures.
All institutional examples below have stipulated rules and evidence;
they state no current rule of law.
The set-sized strategy-existence results use ordinary classical
set theory with choice. The finite examples do not need a choice
principle to select their displayed moves.

\section{Argument construction and the dialectical interface}
\label{sec:p18-lis-interface}

\begin{hbdefinition}{def:p18-lis-argument-interface}{Language, arguments, attacks, and defeat}
Let \(\mathcal L\) be a propositional language with classical negation.
For a literal \(l\), write \(\overline l\) for its complementary
literal; this notation cancels a double negation.
Rules have distinct names and finite lists of premises.
A strict rule \(r:\phi_1,\ldots,\phi_k\to\psi\) is an admitted
truth-preserving inference; a presumptive rule
\(r:\phi_1,\ldots,\phi_k\Rightarrow\psi\) is an admitted defeasible
inference. Empty lists are allowed.
Which rules are admitted and which premises are available are inputs.

An argument is a finite rooted, ordered inference tree whose leaves
are available premises or empty-premise rule applications and whose
internal nodes instantiate admitted rules. Its conclusion
\(\operatorname{Conc}(A)\) is its root formula.
Premise occurrences are marked ordinary, assumption, or certain.
An assumption is admitted defeasibly in the absence of a specified
counterconsideration. A certain premise is protected by the selected
procedure. Neither marking asserts an empirical truth.

Subarguments are subtree \emph{occurrences}, identified by their
addresses in the tree. Thus repeated copies of an identical derivation
have separate occurrences with unique parents.
Write \(\operatorname{Sub}(A)\) for these occurrences, including the root.
The argument universe \(\mathcal A\) and its subargument occurrences
are supplied with this structure. Arguments available in a particular
game are fixed throughout that game.

An attack record identifies an attacking argument, a target
subargument occurrence, and its type:
\begin{enumerate}
\item a contrary to an ordinary premise or assumption;
\item a rebuttal of a presumptive conclusion, including an intermediate
one;
\item an undercutter of a named presumptive application.
\end{enumerate}
An undercutter can use a distinct formula
\(\neg\operatorname{app}(r,\iota)\), where \(\iota\) identifies the
application. It attacks applicability, not the classical implication
asserted by a strict inference.
Certain premises are not attackable where the selected procedure
protects them.

The adjudicated defeat relation \(D\subseteq\mathcal A^2\) is an
additional input recording successful attacks. Write
\[
C\mathrel{D_s}A\quad\Longleftrightarrow\quad
C\mathrel D A\ \text{and not }A\mathrel D C .
\]
Thus strict defeat means defeat without reverse defeat, not merely
that the attacker used a strict inference.
An unresolved comparison may produce mutual defeat.
The comparative grounds may concern rule priority, evidence strength,
or another stated criterion. They are not determined by LIS itself.

Defeat at a subargument propagates to the superargument containing
that occurrence. For each permitted pair \(CDA\), supply the unique
smallest attacked occurrence \(t(C,A)\) in \(A\).
This is the 2001 presentation's single-target restriction.
Pairs with two incomparable minimal targets require a different
extension and are outside this interface.
Whenever a subargument of \(C\) also defeats \(A\), it must identify
the same target occurrence as \(C\).
\end{hbdefinition}

This interface permits an explicit comparison matrix in a worked case.
It does not silently choose the last-link or weakest-link preferences
of the separate ASPIC construction.
The 2006 IS summary uses extended logic-program rules with strong
negation and weak premises \(\mathord{\sim}l\), where an argument
for \(l\) attacks the weak assumption.
It also distinguishes inapplicability and rebuttal, with rule priorities.
Its original systems permit arguments about priorities themselves.
Their source rule construction is supplied in
Chapter~\ref{sec:p18-is1996}, followed by the evolving comparisons
in Chapter~\ref{sec:p18-is-dynamic} and the corrected dialogue
in Chapter~\ref{sec:p18-is-dynamic-completeness}.
These mechanisms are not obtained merely by entering a fixed
matrix \(D\).

The 2009 rule-exception presentation can attach an assumption
\(\neg\operatorname{Exc}(r)\) to a rule and use a further argument
for \(\operatorname{Exc}(r)\) to block it.
The 2001 examples also use an ordinary negative exception premise,
supported if necessary by an empty-premise presumptive argument.
These representations differ in how counterevidence operates.
Specify the representation and its attack record; do not replace an
ordinary premise with a weak assumption without changing the game data.

\section{The fixed-role game and grounded justification}
\label{sec:p18-lis-base-game}

\begin{hbdefinition}{def:p18-lis-base-game}{The corrected fixed-role argument game}
A dialogue starts with \(P\) moving an argument \(A\), and then
alternates \(O,P,O,\ldots\).
Every later move attacks the immediately preceding argument:
\(O\) may use any defeater, while \(P\) must use a strict defeater.
\(P\) may not repeat an argument previously moved by \(P\).
\(O\) may repeat its arguments.
All moves use the fixed universe and relation.

A finite dialogue ends when the player whose turn it is has no
legal move; the preceding player wins.
An infinite dialogue is not a terminal win for either player.
A strategy selects a legal move at each of its player's decision
histories. It is winning if every maximal play compatible with it
ends with the other player unable to move.
Winning one particular play is weaker than having a winning strategy.
\end{hbdefinition}

For comparison, define the defense operator on subsets of
\(\mathcal A\) by
\[
F(S)=\{A:\ \text{for every }BDA
             \text{ there is }C\in S\text{ with }CDB\}.
\]
Its least fixed point \(J\) is the grounded extension.
One construction starts at \(J_0=\varnothing\), uses
\(J_{\alpha+1}=F(J_\alpha)\), and takes unions at limit ordinals.
It stabilizes because \(\mathcal A\) is a set.
The first entry stage of an element of \(J\) is its rank.

\begin{hbtheorem}{thm:p18-lis-grounded-game}{proved=annex:p18.lis-grounded-game}{Fixed-role game correspondence}
For an arbitrary set-sized defeat graph, \(A\in J\) if and only if
\(P\) has a winning strategy in the corrected fixed-role game
starting with \(A\). The assertion does not require finite branching
or a finite argument universe.
\end{hbtheorem}

\begin{hbtheorem}{thm:p18-lis-opponent-repetition}{proved=annex:p18.lis-opponent-repetition}{Why opponent repetition must remain available}
If the fixed-role game instead forbids each player to repeat its own
arguments, it can declare a proponent winner for an argument outside
the grounded extension. This occurs already in a five-argument graph.
\end{hbtheorem}

The fixed-role statuses are:
\emph{justified} for \(A\in J\);
\emph{overruled} for a nonjustified \(A\) defeated by a member of \(J\);
and \emph{defensible} for the remaining arguments.
A formula is justified when some argument for it is justified.
If there is no such argument but a defensible one, call it defensible.
If it has at least one argument and all its arguments are overruled,
call it overruled. If it has no argument, record it as unsupported.
The last convention makes the empty-family case explicit.
Being overruled does not itself establish the formula's negation:
an undercutter can defeat its only argument without supporting a
contrary conclusion.

\section{Issue-relative allocation and historical roles}
\label{sec:p18-lis-burden-input}

\begin{hbdefinition}{def:p18-lis-burden-profile}{The corrected burden-profile domain}
Let \(Q\subseteq\mathcal L\) be a specified set of issues, closed
under complementation and containing no classical tautology.
Here complementation means classical negation with adjacent
double negations cancelled; use this normalization also in the
separation condition below.
For \(S\subseteq Q\), define
\[
\operatorname{Cn}_Q(S)=
 \{\phi\in Q:S\models_{\rm CL}\phi\}.
\]
The parties are \(\pi\) and \(\delta\); write \(\bar p\) for
the other party.
A history is a finite sequence of pairs \((p,A)\), including the
empty history \(\epsilon\).
A burden profile is a specified function
\[
B:\operatorname{Hist}\times\{\pi,\delta\}\longrightarrow\mathcal P(Q).
\]
For every history \(h\), each \(B(h,p)\) is classically consistent
and \(\operatorname{Cn}_Q\)-closed.
If \(\phi\in B(h,p)\), neither \(\phi\) nor \(\neg\phi\) belongs
to \(B(h,\bar p)\).

Retain the source's two interstage constraints.
If \(h'=h\cdot(p,A)\), a burden of the other party cannot disappear:
\(B(h,\bar p)\subseteq B(h',\bar p)\).
A burden \(\phi\in B(h,p)\) of the mover can disappear only if
some subargument of the moved \(A\) concludes \(\phi\).
New burdens can be added subject to consistency and separation.
These constraints regulate a supplied allocation; they do not
derive an institution's reasons for allocating it.
\end{hbdefinition}

\begin{hbtheorem}{thm:p18-lis-burden-domain}{proved=annex:p18.lis-burden-domain}{Closure obstruction and a nonempty corrected domain}
Two burden sets that are both closed under unrestricted classical
consequence cannot satisfy the prohibition on shared burdens:
each would contain every tautology.
In contrast, if \(Q\) consists of complementary literals, every
consistent subset of \(Q\) is \(\operatorname{Cn}_Q\)-closed.
Constant separated consistent literal profiles therefore satisfy
the corrected conditions and their retention constraints.
\end{hbtheorem}

The issue restriction is an additional, explicit domain premise.
It preserves the source's intended separation and its examples of
particular allocated issues while avoiding a demand that every logical
truth become a burden of both parties.
It does not establish the existence of an allocation satisfying
the unrepaired unrestricted-classical condition.
The input profile can still depend on the entire history; it is
not required to be constant, finite-state, or computable.

\begin{hbdefinition}{def:p18-lis-roles}{Roles of subtree occurrences when a move is made}
For a candidate move \(M_i=(p_i,A_i)\), let \(h_{i-1}\) be its
preceding history. Assign a role
\(\rho_i(U)\in\{P,O\}\) to every occurrence \(U\in\operatorname{Sub}(A_i)\).
Roles of earlier moves are retained as assigned when those moves
were made.
\begin{enumerate}
\item At \(i=1\), every occurrence has role \(P\).
\item At \(i>1\), if
\(\operatorname{Conc}(U)\in B(h_{i-1},p_i)\), assign \(P\).
\item Otherwise, if the subtree argument at \(U\) defeats \(A_{i-1}\),
assign the opposite of
\(\rho_{i-1}(t(A_i,A_{i-1}))\).
\item Otherwise assign the role of \(U\)'s immediate parent occurrence
inside the current argument \(A_i\).
\end{enumerate}
For a noninitial candidate, require \(A_iDA_{i-1}\) before applying
this definition. Thus the root is covered by the second or third
case; the fourth case is used only for proper subarguments.
An explicit burden takes precedence over role inheritance.
\end{hbdefinition}

\begin{hbtheorem}{thm:p18-lis-role-totality}{proved=annex:p18.lis-role-totality}{Roles are well defined on the declared target domain}
For every finite candidate history meeting the defeat prerequisites,
the role definition assigns exactly one role to each subtree occurrence.
The assignment can be constructed in increasing move order and,
within each move, from root toward leaves.
Without a unique minimal target, the source's rule can be ambiguous.
\end{hbtheorem}

\section{All litigation moves, closure, and winning}
\label{sec:p18-lis-game}

\begin{hbdefinition}{def:p18-lis-game}{The changing-role litigation game}
The first move is \((\pi,A_1)\). Physical parties then alternate;
their argument roles need not alternate.
A candidate move \(M_i=(p_i,A_i)\), \(i>1\), is legal exactly when:
\begin{enumerate}
\item \(p_i\) is the party whose turn it is and \(A_i\in\mathcal A\);
\item \(A_iDA_{i-1}\);
\item if the root role \(\rho_i(A_i)\) is \(P\), then
\(A_iD_sA_{i-1}\);
\item there is no earlier move \(M_j=(p_i,A_i)\), \(j<i\),
whose root role at that earlier move was \(P\).
\end{enumerate}
The last condition allows \(O\)-then-\(P\) and \(O\)-then-\(O\)
repetitions, while prohibiting any same-party reuse after that
argument has been moved in role \(P\).
It does not confuse physical party with dialectical role.

The complete state is the finite history together with its recursively
computed role annotations. The supplied \(B\), argument data,
and \(D\) determine its possible successors.
All legal continuations form the game arena.
There are no additional concession, withdrawal, evidence-introduction,
or judicial moves inside this dialectical game.
Those belong to an outer procedure supplying its input states.

Terminal and winning conditions are those of the fixed-role game:
a party wins a finite maximal play when the other party has no legal
continuation; an infinite play is not a terminal win.
Say that \(\pi\) \emph{establishes \(A_1\) in this litigation game}
when it has a winning strategy from that initial move.
Do not identify this status with grounded membership without a
separate correspondence.
\end{hbdefinition}

A game tree records histories, not just argument names.
Identical arguments reached with different burdens, previous
proponent uses, or role annotations can have different continuations.
A proof tree inside one argument is a different object again.

\begin{hbdefinition}{def:p18-lis-attractor}{The exact winning-set construction}
Fix a legal initial move and let \(H\) be all reachable finite
histories. Write \(\operatorname{Succ}(h)\) for its legal successors
and \(\operatorname{turn}(h)\) for the next physical party.
For a party \(p\), let
\[
T_p=\{h:\operatorname{Succ}(h)=\varnothing,\
                   \operatorname{turn}(h)=\bar p\}.
\]
Define the monotone operator on \(\mathcal P(H)\):
\[
\begin{split}
\Phi_p(X)=T_p
&\cup\{h:\operatorname{turn}(h)=p,\
                 \operatorname{Succ}(h)\cap X\ne\varnothing\}\\
&\cup\{h:\operatorname{turn}(h)=\bar p,\
     \varnothing\ne\operatorname{Succ}(h)\subseteq X\}.
\end{split}
\]
Starting with \(X_0=\varnothing\), iterate \(\Phi_p\) at successors
and take unions at limits. Let \(W_p\) be the stabilized set.
The nonempty-successor condition prevents a losing terminal from
being accepted by vacuous universal quantification.
\end{hbdefinition}

\begin{hbtheorem}{thm:p18-lis-winning-construction}{proved=annex:p18.lis-winning-construction}{Complete winning characterization}
For the fully declared, possibly infinite history arena, \(h\in W_p\)
if and only if \(p\) has a strategy forcing a finite terminal win
from \(h\). The sets \(W_\pi,W_\delta\) are disjoint.
Histories outside both sets need not terminate and are not assigned
to a winner by this construction.
If a finite arena with \(N\) states is explicitly supplied, iteration
stabilizes after at most \(N\) strict increases.
This is not a claim that arbitrary history-dependent profiles produce
a finite arena.
\end{hbtheorem}

This proof supplies a complete independent evaluation construction.
The source defines litigation winning dialectically and expressly
does not establish a general declarative equivalence to grounded
semantics. The construction above does not attribute an unproved
equivalence to it.

\begin{hbtheorem}{thm:p18-lis-grounded-separation}{proved=annex:p18.lis-grounded-separation}{A litigation win can differ from grounded justification}
There is a three-argument, single-target instance with constant
consistent literal burdens in which \(\pi\) has a winning litigation
strategy for \(A\), while \(A\) is not grounded-justified.
\end{hbtheorem}

\section{Complete worked games}
\label{sec:p18-lis-worked}

\subsection{Tied evidence and an issue-specific change of role}
\label{sec:p18-lis-tied-evidence}

Use three empty-premise presumptive arguments:
\(A\) concludes \(\neg b\), \(B\) concludes \(b\), and \(C\)
concludes \(\neg b\).
The complete defeat relation is
\[
BDA,\qquad BDC,\qquad CDB,
\]
with no other pair. Thus \(B\) strictly defeats \(A\), while
\(B,C\) defeat each other.
This can represent a comparison in which \(B\)'s support outranks
\(A\)'s, but \(B\)'s and \(C\)'s support remain tied.
Every target is the root.
Let \(Q=\{b,\neg b\}\), \(B(h,\delta)=\{b\}\), and
\(B(h,\pi)=\varnothing\) at every history.

The complete play is:
\[
(\pi,A)_P,\quad(\delta,B)_P,\quad(\pi,C)_O .
\]
The subscript records the role, not the party.
\(A\) is initially \(P\); \(B\) is \(P\) because \(\delta\)
owns \(b\), so its strict defeat is required and available.
\(C\) has no explicit burden and attacks a \(P\)-role root,
so it is \(O\) and mutual defeat suffices.
At the next turn the only possible attacking argument is \(B\).
The same party \(\delta\) already used \(B\) in role \(P\),
so it cannot repeat it. There is no other legal move:
\(\pi\) wins. Since these responses are the only ones, the play
also specifies a winning strategy, not merely a fortunate branch.
The grounded calculation differs, as proved in the annex.

\subsection{Nested exception and document issues}
\label{sec:p18-lis-document-issues}

In a fictional eligibility dispute, \(A\) supports an award from
an application and the ordinary default-supported premise
\(\neg e\), absence of an exception.
\(B\) supports \(e\) from a document \(d\);
\(C\) supports \(\neg d\) from a recorded defect in its stamp;
\(D\) supports \(d\) from an independent authentication record.
The tree occurrences are:
\[
A=[a,\ [\Rightarrow\neg e]\Rightarrow w],\quad
B=[[r_d\Rightarrow d]\Rightarrow e],\quad
C=[r_c\Rightarrow\neg d],\quad D=[r_D\Rightarrow d].
\]
The displayed notation abbreviates distinct named rules and their
indicated premise arguments; the names \(r_d,r_c,r_D\) here
denote available evidence contents, not unproved legal propositions.
Include all subarguments. Evidence leaves and the application
premise have no attackers.

Supply this complete pattern of successful defeats:
\(B\) defeats the \(\neg e\) subargument and hence \(A\);
\(C\) defeats \(B\) at its \(d\) occurrence and hence defeats
that subargument;
\(C,D\) mutually defeat each other.
\(BDA\) and \(CDB\) are strict.
No other pair succeeds. In particular the weaker document
subargument of \(B\) does not defeat \(C\).
This states the evidential comparison instead of presuming that
every document has equal weight.

Use constant issue burdens \(\{w,\neg d\}\) for \(\pi\) and
\(\{e\}\) for \(\delta\), with \(Q\) the complementary literals
for \(w,d,e\).
They are consistent and separated.
The only main-line play is
\[
(\pi,A)_P,\quad(\delta,B)_P,\quad
(\pi,C)_P,\quad(\delta,D)_O .
\]
The first role is fixed by the opening convention.
\(B\)'s exception conclusion has \(\delta\)'s burden.
\(C\)'s document challenge has \(\pi\)'s burden.
\(D\) has no allocated burden and attacks the \(P\)-role \(C\),
so it is \(O\).
Its mutual defeat is sufficient.
The only response attacking \(D\) is \(C\), whose same-party
reuse after a \(P\) move is forbidden. Thus \(\delta\) wins.

Change only the comparison so that \(C\) strictly defeats \(D\).
Then \(D\) no longer defeats \(C\), and the play terminates after
\(C\) with \(\pi\) winning. This comparison change differs from
reassigning either party's burden.
Alternatively, remove \(\neg d\) from \(\pi\)'s initial burden:
then \(C\), responding to \(B\)'s \(P\) target, has role \(O\).
Its required force changes even though the same physical party
and evidence are involved. In that alternative profile, \(D\)
would respond to an \(O\)-role target and itself have role \(P\).
Its merely mutual defeat of \(C\) is then insufficient.
The game ends after \(C\), with \(\pi\) winning.

\subsection{Where a subargument is attacked matters}
\label{sec:p18-lis-target-specific}

Let \(A=[A_o,A_x\Rightarrow c]\), with
\(A_o=[p\Rightarrow o]\) and \(A_x=[v\Rightarrow x]\).
The contents represent possession, a claimed entitlement,
a witness report, alleged damage, and a claim for compensation
in a fictional rule system.
For this case take \(Q\) to contain both polarities of
\(p,o,v,x,c\).
Initially assign \(\{o,x,c\}\) to \(\pi\) and nothing to \(\delta\).
After \(\pi\) moves \(A\), a stipulated outer allocation removes
\(o\) from \(\pi\) and assigns \(\neg o\) to \(\delta\).
It retains \(x,c\) for \(\pi\).
Removal is permitted by the profile constraint because \(A_o\)
concluded \(o\). All earlier role annotations remain unchanged.
Retain this new allocation at later histories in this example.

A response concluding \(\neg o\) has role \(P\) because of
\(\delta\)'s new explicit burden and must strictly defeat its target.
A response concluding \(\neg x\) has no such burden and targets
the originally \(P\)-role \(A_x\); it has role \(O\), so defeat
suffices. A response to the witness premise \(v\) likewise receives
the opposite of that occurrence's inherited \(P\) role.
Thus different responses to the same moved argument can require
different strengths. The outer allocation is stipulated here;
LIS does not infer its institutional justification from possession.

For complete branch outcomes, add three empty-premise candidate
counterarguments \(N_o,N_x,N_v\), concluding \(\neg o,\neg x,\neg v\).
Let the complete defeat matrix consist of mutual pairs
\[
(N_o,A_o),\ (N_o,A),\quad
(N_x,A_x),\ (N_x,A),\quad
(N_v,v),\ (N_v,A_x),\ (N_v,A),
\]
where each listed pair includes both directions and \(v\)
also denotes the premise argument for the witness report.
There are no other successful defeats.
These pairs explicitly include propagation to superarguments.
At \(\delta\)'s first turn, \(N_o\) is inadmissible because its
role is \(P\) but its defeat is only mutual.
Either \(N_x\) or \(N_v\) is legal in role \(O\).
After either choice, every available defeater of that move is
only mutual, while the required response has role \(P\):
for \(A_x\) its burden \(x\) gives \(P\), and for the witness
argument its attack on the preceding \(O\)-role target gives \(P\).
The full argument \(A\) is also barred by its earlier \(P\) use.
No strict legal reply exists, so either legal branch ends with
\(\delta\) winning. If \(N_o\) alone were available, it would
remain illegal and \(\pi\) would win immediately after \(A\).

\section{Production, persuasion, and tactical status at stages}
\label{sec:p18-lis-stage-burdens}

The 2009 stage account fixes the information pool while evaluating
one game, then permits the outer proceeding to produce a new pool.
Let a stage record supply the available arguments, their attack
and defeat data, and the applicable burden profile.
At the final stage, the designated game evaluates the claim.
An intermediate tactical assessment evaluates what would happen
if that stage's information and predicted comparative standards
were final. It is conditional on that forecast.

For a persuasion burden, record the issue, party, and applicable
standard. In the fixed-role account it requires a justified
argument at the final stage. With distributed LIS roles, use the
specified litigation winning condition, preserving its distinction
from grounded membership.
The standard affects adjudication of conflicting evidence:
a slight advantage need not yield strict defeat under every
standard. No universal numerical conversion is supplied.

A production burden is a threshold for admitting an evidential
argument as internally adequate for consideration.
Its substantive threshold is an input, not a theorem of the game.
The 2009 treatment models a negative ruling by a strictly defeating
argument that the selected procedure protects from further challenge.
Once the threshold is met, subsequent adverse evidence affects
the merits or tactical forecast; it does not erase the earlier
production event.

\begin{hbtheorem}{thm:p18-lis-protected-production}{proved=annex:p18.lis-protected-production}{A protected production ruling settles its target}
If an available argument \(U\) has no defeaters and strictly defeats
an evidential argument \(E\), then \(U\) is grounded-justified and
\(E\) is overruled. In either argument game, a legal reply with \(U\)
ends that branch in favor of its mover.
\end{hbtheorem}

\begin{hbexample}{def:p18-lis-stage-example}{A complete four-stage evidential sequence}
Use the fixed-role case in which one party owns the main persuasion
burden. Let \(M\) support its claim; \(E\) gives an exception;
\(R\) counters that exception; and \(U\) undercuts \(R\).
All targets are unique and the full successful defeat chain is
\[
U D_s R D_s E D_s M .
\]
No additional defeats occur. At stage 0 only \(M\) is available:
its proponent wins immediately.
At stage 1, \(E\) is added and passes the stipulated production
threshold. The play \(M,E\) ends against the main claimant.
At stage 2, \(R\) is added; the complete play \(M,E,R\) ends
for that claimant.
At stage 3, \(U\) is added; \(M,E,R,U\) ends against it.
There is exactly one continuation at every turn, so these are
complete games and not selected favorable branches.

The exception's production threshold was passed at stage 1
and stays passed. The original persuasion allocation is unchanged.
The need for further evidence alternates with the stage forecasts.
If instead the production ruling is an undefeated argument
against \(E\), the protected-ruling theorem describes the
procedural rejection branch.
\end{hbexample}

An overruled main claim in this sequence is not proof of its
classical negation. In particular, \(U\)'s conclusion may only
deny the applicability of \(R\)'s inference. The assessment
concerns what the available arguments establish under the
declared procedure.

\section{Scope and the linked original-source constructions}

The 2001 article's changing allocations concern a supplied
history-dependent profile. The 2009 distinction between fixed
persuasion/production responsibilities and changing tactical
pressure concerns a specified proceeding. They are compatible
only after stating which changes the outer procedure authorizes.
One must not describe every tactical reversal as a reassignment
of a persuasion burden.

The 2009 account additionally represents uncontested claims as
protected premises where a procedure gives silence that effect;
excludes unclaimed issues where a burden of claiming requires
timely assertion; and prevents a bare premise from satisfying
a burden of argument that requires a derivation.
These are different outer procedures and need explicit choices.
The game supplies none of them automatically.

\begin{proofboundary}{def:p18-lis-next-boundary}{Dialectical interfaces and remaining historical obligations}
The argument interface and corrected issue-relative game above
have complete definitions and proofs for their declared domain.
They do not assert a general legal allocation rule or a
correspondence to the separate ASPIC or DDF systems.

The rule-generated fixed-priority IS construction and its exact
fixed-role bridge are given in Chapter~\ref{sec:p18-is1996}.
Chapter~\ref{sec:p18-is-dynamic} supplies dynamic comparisons,
their semantics, applicability alternatives and their limitations.
Theorem~\ref{thm:p18-is-compression-countermodel} refutes the
original single-response completeness claim on its full domain;
Theorem~\ref{thm:p18-is-licensing-complete} proves the corrected
sequential dialogue correspondence. The changing-role LIS
interface still requires its own designated targets and burden
data; these are not supplied by the IS bridge automatically.

The cited Dung dialog-game report retains its separate exact
1993/1994 version question and historical proof-reading obligation.
The full 2006 presumptions/backing presentation, sensitive direct
exclusion redesign, and broader implementation and stable-semantics
comparisons remain distinct source developments.
The independent grounded-game proof here establishes its stated
mathematics without claiming to reproduce that unread Dung proof.
\end{proofboundary}
